\chapter*{論文要旨}
\addcontentsline{toc}{chapter}{論文要旨}

偉人の伝記というと，ナポレオンとかアレキサンドロスとか，グラッドストーンというようなのばかりで，学者のはほとんど無いと言ってよい．
なるほどナポレオンやアレキサンドロスのは，雄であり，壮である．
しかし，いつの世にでもナポレオンが出たり，アレキサンドロスの出ずることは出来ない．
文化の進まざる時代の物語りとして読むには適していても，修養の料にはならない．
グラッドストーンのごときといえども，一国について見れば二，三人あり得るのみで，しかも大宰相たるは一時に一人のみしか存在を許さない．
これに反して，科学者や哲学者や芸術家や宗教家は，一時代に十人でも二十人でも存在するを得，また多く存在するほど文化は進む．
ことに科学においては，言葉を用うること少なきため，他に比して著しく世界的に共通で，日本での発見はそのまま世界の発見であり，詩や歌のごとく，外国語に訳するの要もない．

これらの理由により，科学者たらんとする者のために，大科学者の伝記があって欲しい．
しかし，科学者の伝記を書くということは，随分むずかしい．というのは，まず科学そのものを味った人であることが必要であると同時に多少文才のあることを要する．
悲しいかな，著者は自ら顧みて，決してこの二つの条件を備えておるとは思わない．
ただ最初の試みをするのみである．
　
科学者の中で，特にファラデーを選んだ理由は，第一に彼は大学教育を受けなかった人で，全くの丁稚小僧から成り上ったのだ．
学界にでは家柄とか情実とかいうものの力によることがない，腕一本でやれるということが明かになると思う．また立身伝ともいえる．
次に彼の製本した本も，筆記した手帳も，実験室での日記も，発見の時に用いた機械も，それから少し変ってはいるが，実験室も今日そのまま残っている．
シェーキスピアやカーライルの家は残っている．
ゲーテ，シルレルの家もあり，死んだ床も，薬を飲んだ杯までもある（真偽は知らないが）．ファラデーのも，これに比較できる位のものはある．
科学者でファラデーほど遺物のあるのは，他に無いと言ってよい．
それゆえ，伝記を書くにも精密に書ける．
諸君がロンドンに行かるる機会があったら，これらの遺物を実際に見らるることも出来る．
　
第三に，何にか発見でもすると，その道行きは止めにして，出来上っただけを発表する人が多い．
感服に値いしないことはないが，これでは，後学者が発見に至るまでの着想や推理や実験の順序方法について，貴ぶべき示唆を受けることは出来ない．
あたかも雲に聳（そび）ゆる高塔を仰いで，その偉観に感激せずにはいられないとしても，さて，どういう足場を組んで，そんな高いものを建て得たかが，判らないのと同じである．

ファラデーの論文には，いかに考え，いかに実験して，それでは結果が出なくて，しまいにかくやって発見した，というのが，偽らずに全部書いてある．これでこそ発見の手本にもなる．
またファラデーの伝記は決して無味乾燥ではない．電磁気廻転を発見して，踊り喜び，義弟をつれて曲馬見物に行き，入口の所でこみ合って喧嘩をやりかけた壮年の元気は中々さかんである．
莫大の内職をすて，［＃「莫大の内職をすて，」は底本では「莫大の内職をすて，」］宴会はもとより学会にも出ないで，専心研究に従事した時代は感嘆するの外はない，晩年に感覚も鈍り，ぼんやりと椅子（いす）にかかりて，西向きの室から外を眺めつつ日を暮らし，終に眠るがごとくにこの世を去り，静かに墓地に葬られた頃になると，落涙を禁じ得ない．

前編に大体の伝記を述べて，後編に研究の梗概（こうがい）を叙することにした．
